\section{Calculo factor $R^*$ para sismo en X e Y}
Con el valor calculado en el apartado anterior, se puede determinar el factor de reducción de la fuerza sísmica, $R^*$, utilizando la fórmula (6-10) de la norma NCh 433 \cite{NCh433:1996}:

\begin{equation}
    R^* = \frac{T*}{0.1 \cdot T_0 + \frac{T*}{R_0}}
\end{equation}

Donde $T*$ es el período del modo con mayor masa traslaciona equivalente de análisis y $T_0$ es un parámetro de la norma que esta asociado con la ubicación del edificio.

Para poder determinar el $R^*$, se corrió el modelo en ETABS y se obtuvieron los siguientes resultados en cuanto a los períodos de los modos de vibración:

\begin{table}[H]
    \centering
    \begin{tabular}{ccc}

        Modo  & Eje traslacional & Período (s) \\
        \hline
        1 & torcional & 1.251 \\
        2 & eje X & 1.14 \\
        3 & eje Y & 0.853 \\
    \end{tabular}
    \caption{Períodos de los modos de vibración obtenidos en ETABS}
    \label{tab:periodos}
\end{table}

Para este análisis, se va a considerar el modo 2 y 3. Para poder calcular el $R^*$ respectivo, se necesita determinar el $T_0$ asociado que se encuentra en la tabla \ref{tab:parametros_sismicos}. Con todos los datos necesarios, se calcula el $R^*$ para cada modo:
\begin{itemize}
    \item Para el modo 2 (eje X):
    \begin{equation}
        R^* = \frac{1.14}{0.1 \cdot 0.15 + \frac{1.14}{11}} = 10.609
    \end{equation}

    \item Para el modo 3 (eje Y):
    \begin{equation}
        R^* = \frac{0.853}{0.1 \cdot 0.15 + \frac{0.853}{11}} = 10.217
    \end{equation}
\end{itemize}

Sin embargo, para el análisis sísmico en ETABS, se ocupa un el inverso del $R^*$, por lo tanto, el valor a ocupar en el análisis sísmico es:
\begin{equation}
    \frac{1}{R^*} = \frac{1}{10.609} = 0.0943
\end{equation}
\begin{equation}
    \frac{1}{R^*} = \frac{1}{10.217} = 0.0979
\end{equation}

Estos datos se ingresan en los espectros respectivos para cada eje y se obtiene el corte basal en cada uno de ellos. Estos resultados se muestran a continuación:

\begin{table}[H]
    \centering
    \begin{tabular}{ccc}
        Eje & Corte basal (ton) \\
        \hline
        X & 160.34 \\
        Y & 201.92 \\
    \end{tabular}
    \caption{Corte basal obtenido para cada eje}
    \label{tab:corte_basal}
\end{table}

Con estos valores de corte basal, 